% ═══════════════════════════════════════════════════════════════
% GREENHOUSE ARBEITSBLATT TEMPLATE
% Für Stundenleistungen, Leistungskontrollen, Übungsblätter
% ═══════════════════════════════════════════════════════════════

\documentclass[11pt, a4paper]{article}

% ═══════════════════════════════════════════════════════════════
% PAKETE
% ═══════════════════════════════════════════════════════════════

\usepackage[utf8]{inputenc}
\usepackage[T1]{fontenc}
\usepackage[ngerman]{babel}
\usepackage{helvet}                     % Serifenlose Schrift
\renewcommand{\familydefault}{\sfdefault}

\usepackage[margin=2cm]{geometry}
\usepackage{amsmath, amssymb}
\usepackage{siunitx}                    % SI-Einheiten
\sisetup{locale = DE, per-mode = symbol}

\usepackage{graphicx}
\usepackage{tikz}
\usepackage{pgfplots}
\pgfplotsset{compat=1.18}

\usepackage{enumitem}                   % Anpassbare Listen
\usepackage{tabularx}
\usepackage{booktabs}
\usepackage{multirow}

\usepackage{xcolor}
\usepackage{framed}
\usepackage{tcolorbox}

\usepackage{fancyhdr}
\usepackage{lastpage}

% ═══════════════════════════════════════════════════════════════
% FARBEN (sciencesim-konform)
% ═══════════════════════════════════════════════════════════════

\definecolor{physik}{HTML}{2c5aa0}
\definecolor{chemie}{HTML}{2a7a4b}
\definecolor{informatik}{HTML}{b35c00}
\definecolor{grau}{HTML}{888888}
\definecolor{hellgrau}{HTML}{f5f5f5}

% Aktive Fachfarbe setzen (Standard: Physik)
\newcommand{\fachfarbe}{physik}
% Zum Wechseln: \renewcommand{\fachfarbe}{chemie}

% ═══════════════════════════════════════════════════════════════
% VARIABLEN
% ═══════════════════════════════════════════════════════════════

\newcommand{\thema}{Bewegungslehre}
\newcommand{\fach}{Physik}
\newcommand{\klasse}{10}
\newcommand{\typ}{Stundenleistung}          % SL, LK, AB, UE
\newcommand{\dauer}{45 Minuten}
\newcommand{\maxpunkte}{50}
\newcommand{\datum}{\today}

% Notengrenzen (MR 6-Noten-Skala)
\newcommand{\noteEins}{48}
\newcommand{\noteZwei}{40}
\newcommand{\noteDrei}{30}
\newcommand{\noteVier}{20}
\newcommand{\noteFuenf}{10}

% ═══════════════════════════════════════════════════════════════
% KOPF- UND FUSSZEILE
% ═══════════════════════════════════════════════════════════════

\pagestyle{fancy}
\fancyhf{}
\renewcommand{\headrulewidth}{0.4pt}
\renewcommand{\footrulewidth}{0.4pt}

\lhead{\fach{} -- Klasse \klasse}
\chead{\typ}
\rhead{\datum}

\lfoot{Greenhouse Schools Rostock}
\cfoot{}
\rfoot{Seite \thepage\ von \pageref{LastPage}}

% ═══════════════════════════════════════════════════════════════
% BEFEHLE
% ═══════════════════════════════════════════════════════════════

% Punkteplatzhalter am Rand
\newcommand{\punkte}[1]{%
    \hfill\textcolor{grau}{[\hspace{0.5cm}/#1~BE]}%
}

% Sternchen-Differenzierung
\newcommand{\stern}[1]{%
    \textsuperscript{\textcolor{\fachfarbe}{(#1)}}%
}
% Verwendung: \stern{*}, \stern{**}, \stern{***}

% Aufgabenbox
\newenvironment{aufgabe}[2]{%
    \vspace{10pt}
    \noindent\textbf{\textcolor{\fachfarbe}{Aufgabe #1}} #2
    \vspace{5pt}
    \begin{list}{}{\leftmargin=0pt}
    \item[]
}{%
    \end{list}
}

% Teilaufgaben
\newlist{teilaufgaben}{enumerate}{2}
\setlist[teilaufgaben,1]{label=\alph*), leftmargin=2em}
\setlist[teilaufgaben,2]{label=\roman*), leftmargin=2em}

% Leerzeilen für Antworten
\newcommand{\antwortzeilen}[1]{%
    \vspace{2pt}
    \foreach \n in {1,...,#1}{%
        \noindent\rule{\textwidth}{0.4pt}\vspace{8pt}
    }
}

% Formelbox
\newcommand{\formelbox}[1]{%
    \begin{center}
    \fcolorbox{\fachfarbe}{hellgrau}{\parbox{0.6\textwidth}{\centering #1}}
    \end{center}
}

% Hinweisbox
\newtcolorbox{hinweis}{
    colback=hellgrau,
    colframe=\fachfarbe,
    boxrule=0.5pt,
    arc=0pt,
    left=5pt,
    right=5pt,
    top=5pt,
    bottom=5pt,
    fonttitle=\bfseries,
    title=Hinweis
}

% Operator hervorheben
\newcommand{\operator}[1]{\textbf{\textcolor{\fachfarbe}{#1}}}

% ═══════════════════════════════════════════════════════════════
% DOKUMENT
% ═══════════════════════════════════════════════════════════════

\begin{document}

% ═══════════════════════════════════════════════════════════════
% KOPFBEREICH
% ═══════════════════════════════════════════════════════════════

\begin{center}
    {\Large\bfseries \typ: \thema}
\end{center}

\vspace{5pt}

\noindent
\begin{tabularx}{\textwidth}{|X|X|X|}
\hline
\textbf{Name:} \rule{4cm}{0.4pt} & 
\textbf{Klasse:} \klasse & 
\textbf{Sitzplatz:} P\rule{1cm}{0.4pt} \\
\hline
\end{tabularx}

\vspace{10pt}

% Notenspiegel
\noindent
\begin{tabularx}{\textwidth}{|c|c|c|c|c|c|c|X|}
\hline
\textbf{Note} & 1 & 2 & 3 & 4 & 5 & 6 & \multirow{2}{*}{\textbf{Erreicht:} \rule{1cm}{0.4pt} / \maxpunkte{} BE} \\
\cline{1-7}
\textbf{ab BE} & \noteEins & \noteZwei & \noteDrei & \noteVier & \noteFuenf & 0 & \\
\hline
\end{tabularx}

\vspace{5pt}

\noindent\textcolor{grau}{\small Bearbeitungszeit: \dauer{} | Erlaubte Hilfsmittel: Taschenrechner, Formelsammlung}

\vspace{10pt}
\hrule
\vspace{15pt}

% ═══════════════════════════════════════════════════════════════
% AUFGABEN
% ═══════════════════════════════════════════════════════════════

% --- AUFGABE 1: Multiple Choice ---
\begin{aufgabe}{1}{\stern{*}}
\operator{Kreuze} die richtige Antwort an. \punkte{2}
\end{aufgabe}

\noindent Welche Einheit hat die Geschwindigkeit?

\begin{itemize}[label=$\square$, leftmargin=2em]
    \item $\si{m}$
    \item $\si{m/s}$
    \item $\si{m/s^2}$
    \item $\si{kg \cdot m/s}$
\end{itemize}

\vspace{15pt}

% --- AUFGABE 2: Lückentext ---
\begin{aufgabe}{2}{\stern{*}}
\operator{Ergänze} die Lücken mit den passenden Begriffen. \punkte{4}
\end{aufgabe}

\noindent\textit{Wortliste: Geschwindigkeit, Beschleunigung, Weg, Zeit}

\vspace{5pt}

\noindent 
Bei einer gleichförmigen Bewegung bleibt die \rule{3cm}{0.4pt} konstant.
Der zurückgelegte \rule{3cm}{0.4pt} ist proportional zur \rule{3cm}{0.4pt}.

\vspace{15pt}

% --- AUFGABE 3: Berechnung ---
\begin{aufgabe}{3}{\stern{**}}
\operator{Berechne} die Geschwindigkeit. \punkte{4}
\end{aufgabe}

\noindent Ein Auto fährt eine Strecke von \SI{150}{km} in \SI{2}{h}.

\begin{teilaufgaben}
    \item Gib die gegebenen Größen an.
    \item Berechne die Geschwindigkeit in $\si{km/h}$.
    \item Rechne das Ergebnis in $\si{m/s}$ um.
\end{teilaufgaben}

\antwortzeilen{8}

\vspace{15pt}

% --- AUFGABE 4: Diagramm ---
\begin{aufgabe}{4}{\stern{**}}
\operator{Interpretiere} das Weg-Zeit-Diagramm. \punkte{6}
\end{aufgabe}

\begin{center}
\begin{tikzpicture}
\begin{axis}[
    width=10cm, height=6cm,
    xlabel={$t$ in s},
    ylabel={$s$ in m},
    xmin=0, xmax=10,
    ymin=0, ymax=50,
    grid=major,
    axis lines=left,
    xtick={0,2,4,6,8,10},
    ytick={0,10,20,30,40,50}
]
\addplot[thick, \fachfarbe, mark=*] coordinates {
    (0,0) (2,20) (4,20) (6,20) (8,40) (10,50)
};
\end{axis}
\end{tikzpicture}
\end{center}

\begin{teilaufgaben}
    \item Beschreibe die Bewegung in den Abschnitten I (0--2\,s), II (2--6\,s), III (6--10\,s).
    \item Berechne die Geschwindigkeit im Abschnitt I.
\end{teilaufgaben}

\antwortzeilen{6}

\vspace{15pt}

% --- AUFGABE 5: Transfer ---
\begin{aufgabe}{5}{\stern{***}}
\operator{Analysiere} die Situation und \operator{beurteile} das Ergebnis. \punkte{8}
\end{aufgabe}

\noindent Ein Radfahrer startet mit einer Anfangsgeschwindigkeit von $v_0 = \SI{0}{m/s}$ und beschleunigt gleichmäßig. Nach $t = \SI{10}{s}$ hat er eine Geschwindigkeit von $v = \SI{5}{m/s}$ erreicht.

\begin{teilaufgaben}
    \item Berechne die Beschleunigung.
    \item Berechne den in dieser Zeit zurückgelegten Weg.
    \item Ein PKW benötigt für die gleiche Geschwindigkeitsänderung nur \SI{2}{s}. Vergleiche die Beschleunigungen und beurteile, welche Beschleunigung für den menschlichen Körper angenehmer ist.
\end{teilaufgaben}

\antwortzeilen{10}

% ═══════════════════════════════════════════════════════════════
% FORMELSAMMLUNG (optional auf Rückseite)
% ═══════════════════════════════════════════════════════════════

\newpage

\begin{center}
    {\large\bfseries Formelsammlung}
\end{center}

\vspace{10pt}

\begin{tabularx}{\textwidth}{|l|X|l|}
\hline
\textbf{Größe} & \textbf{Formel} & \textbf{Einheit} \\
\hline
Geschwindigkeit & $v = \dfrac{s}{t}$ oder $v = \dfrac{\Delta s}{\Delta t}$ & $\si{m/s}$ \\[8pt]
\hline
Beschleunigung & $a = \dfrac{\Delta v}{\Delta t}$ & $\si{m/s^2}$ \\[8pt]
\hline
Weg (gleichförmig) & $s = v \cdot t$ & $\si{m}$ \\[8pt]
\hline
Weg (beschleunigt) & $s = \dfrac{1}{2} a t^2 + v_0 t$ & $\si{m}$ \\[8pt]
\hline
Endgeschwindigkeit & $v = a \cdot t + v_0$ & $\si{m/s}$ \\[8pt]
\hline
\end{tabularx}

\vspace{15pt}

\begin{hinweis}
\textbf{Umrechnung:} $\SI{36}{km/h} = \SI{10}{m/s}$ (teilen durch 3,6)

\textbf{Fallbeschleunigung:} $g = \SI{10}{m/s^2}$
\end{hinweis}

\end{document}
